\documentclass[12pt]{amsart}
\pagestyle{empty}
\usepackage[hmargin=1in,vmargin=1in]{geometry}
\usepackage{url,  mathrsfs}
\usepackage{hyperref}
\usepackage{fancyhdr} 
\usepackage[T1]{fontenc}
\usepackage{amssymb}

\usepackage{url,  mathrsfs}
\usepackage{hyperref}
\usepackage{amsmath}
\usepackage{graphicx,  epstopdf}
\usepackage{subfig}
\usepackage{color}
\usepackage{bbm}
\usepackage{subfloat}
\usepackage[draft]{fixme}
\usepackage{bm} 
\usepackage{bbm} 
\usepackage{epigraph}
\usepackage{endnotes}
\usepackage{ifpdf}
\usepackage{array}

\usepackage{multicol}
\usepackage{multirow}
\usepackage{graphicx}
\usepackage{tikz}

\newcommand{\A}{\mathbb{A}}
\newcommand{\Q}{\mathbb{Q}}
\newcommand{\N}{\mathbb{N}}
\newcommand{\Z}{\mathbb{Z}}
\newcommand{\R}{\mathbb{R}}
\newcommand{\F}{\mathbb{F}}
\newcommand{\C}{\mathbb{C}}
\renewcommand{\P}{\mathbb{P}}
\newcommand{\I}{\mathcal{I}}
\newcommand{\V}{\mathcal{V}}
\newcommand{\cH}{\mathcal{H}}
\newcommand{\ol}{\overline}
\newcommand{\J}{\mathcal{J}}
\newcommand{\X}{\mathcal{X}}
\theoremstyle{definition}
\newtheorem{prob}{Problem}
\newtheorem{extraProb}{Extra Problem}

\DeclareMathOperator{\conv}{conv}
\DeclareMathOperator{\cone}{cone}
\DeclareMathOperator{\ext}{ext}




\begin{document}
 \thispagestyle{empty}
\begin{center} {\large \textbf{Math 380 -- Homework  7}} \smallskip \\ 
Due on Thursday, May 28\ \\ \ \\
 \end{center}

You are welcome to talk with other students in the class about problems but should write up solutions on your own. 
Solutions can be handwritten or typed but need to be legible and submitted via Gradescope by the end of the day 
on Thursday. You should justify all your answers in order to receive full credit.  
\begin{flushleft}

{\bf Good practice problems} (not to be turned in): \\
CLO Section 4.3, Exercises 3, 4, 7, 8, 12 \\
CLO Section 4.4, Exercises 1, 4, 11, 14\\
\

% Two ideals I and J of k[x1, . . . , xn] are said to be comaximal if and only if I + J =
% k[x1, . . . , xn].
% a. Show that if k = C, then I and J are comaximal if and only if V(I) ∩ V(J) = ∅. Give
% an example to show that this is false in general.
% b. Show that if I and J are comaximal, then IJ = I ∩ J.
% 198 Chapter 4 The Algebra–Geometry Dictionary
% c. Is the converse to part (b) true? That is, if IJ = I ∩ J, does it necessarily follow that I
% and J are comaximal? Proof or counterexample?
\begin{prob} CLO \S4.3, Exercise 11 parts (a), (b), (c) \smallskip  \\
Two ideals $I, J \subseteq k[x_1,\dots,x_n]$ are \emph{comaximal} if and only if $I + J = k[x_1,\dots,x_n]$.
\begin{enumerate}
    \item[(a)] Over $k = \mathbb C$, $I$ and $J$ are comaximal if and only if $\mathbf V(I)\cap \mathbf V(J) = \varnothing$. Starting from $I$ and $J$ being comaximal, we get $\mathbf V(I) \cap \mathbf V(J) = \mathbf V(I+J) = \varnothing$, since $1 \in I + J = k[x_1,\dots,x_n]$. To prove the other direction, note that by the Weak Nullstellensatz, $\mathbf V(I)\cap \mathbf V(J) = \mathbf V(I+J)=\varnothing$ implies that $1 \in I+J$. Then, $I + J = k[x_1,\dots,x_n]$ since $1\cdot h(x) \in I+J$ for any $h(x)$.

    The Weak Nullstellensatz only works over algebraically closed fields. For a counterexample in $\mathbb R[x,y]$, take $I = \langle x^2+1\rangle$ and $J = \langle y^2+1\rangle$. We have $\mathbf V(I+J) = \varnothing$ since $x^2 +1 \in I+J$ has no solutions, but $1 \not \in I+J$. In fact, $\{ x^2+1,y^2+1\}$ is a Gr\"obner basis under lexicographic order.
    \item[(b)] If $I$ and $J$ are comaximal, then $I \cdot J = I \cap J$.
    \begin{enumerate} 
    \item[$(\subseteq)$] For any element $p \in I\cdot J$, we can write $p = \sum i_k j_k$ where $i_k \in I$ and $j_k \in J$. For all $k$, we get $i_k j_k \in I \cap J$ as ideals absorb multiplication by arbitrary ring elements, and since $I \cap J$ is an ideal closed under addition, $p \in I \cap J$.\item[$(\supseteq)$] Let $a \in I \cap J$, and note that $1\in I + J$ by the fact that $I$ and $J$ are comaximal. Thus, let $i\in I$ and $j \in J$ be such that $i + j = 1$. Then, $a=a\cdot(i+j) = a\cdot i + a\cdot j$. We have that $a \in J$, so $i \cdot a \in I\cdot J$, and similarly $a \in I$ so $a\cdot j \in I \cdot J$. Finally, $a = a\cdot i + a\cdot j \in I \cdot J$.
    \end{enumerate}
    \item[(c)] The converse of (b) is not true. Let $I = \langle 0 \rangle$, and let $J \subsetneq k[x_1,\dots,x_n]$ be any ideal not containing $1$. Then $I \cap J = I\cdot J =\langle 0 \rangle$. However, $1\not\in I+J= J$.
\end{enumerate}
\end{prob}

\bigskip 


% Let V, W ⊆ kn be varieties. Prove that I(V) : I(W) = I(V \ W).
% \[ I:J=\{f\in k[x_1,\dots,x_n] : \forall g \in J, fg \in I\}\]
\begin{prob}  CLO \S4.4, Exercise 8 \smallskip \\
Let $V,W\subseteq k^n$ be varieties. Then $\mathbf I(V):\mathbf I(W) = \mathbf I(V\setminus W)$. 

\begin{enumerate}
    \item[$(\subseteq)$] Let $f \in \mathbf I(V) : \mathbf I(W)$, and let $a \in V \setminus W$. Since $W$ is a variety, $W=\mathbf V(\mathbf I (W))$, and \[a \not \in \mathbf V(\mathbf I (W))  \implies \lnot \forall g \in \mathbf I(W), g(a)=0 \implies \exists g \in \mathbf I(W), g(a)\neq 0.\] Take any such $g \in \mathbf I(W)$. By the definition of $f$, we have $(f\cdot g)(a) = 0$, but since $g(a)\neq 0$, we must have that $f(a)=0$. Thus, $f \in \mathbf I(V\setminus W)$.
    \item[$(\supseteq)$] Let $f \in \mathbf I(V \setminus W)$, and let $g \in \mathbf I(W)$. Let $a \in V$, and first consider $a \not \in W$. In this case, $a \in V \setminus W$ so $f(a) = 0$ and $(f\cdot g)(a)=0$. On the other hand, if $a \in W$, we have $g(a) = 0$, so $(f\cdot g)(a) = 0$. Thus, $f\cdot g \in \mathbf I(V)$, so $f \in \mathbf I(V):\mathbf I(W)$.
\end{enumerate}

% \emph{Note:} You need to work over an arbitrary field $k$, which may not be algebraically closed.\\
\end{prob}

\bigskip 


\end{flushleft}
\end{document}